Dual-Architecture Production-Grade Closed Einstein-Boltzmann Core with Lattice-Resident Torque and Kerr–Gordon Lensing-Photometric Grounding

1. Overall Framework and Dual Architecture

The system is constructed as a strictly dual-layered architecture that unifies a discrete mathematical-numerical “clockwork” with a continuous lensing-photometric “face.”  

Mathematical-Numerical Layer** (the clockwork): exact discrete relations on a tick lattice, phase identities closed to machine precision, and a rapidly screened torque amplitude locked at \(3.170\,\mathrm{km\,s^{-1}\,Mpc^{-1}}\). This layer guarantees internal coherence and the precise local target \(H_{\rm eff}(0)=73.170000\).  
Lensing-Photometric Layer** (the face of the clock): isotropic conformal weight \(W_{\rm iso}\), explicit magnification bias \(\mu_{\rm lens}\), integrated comoving and luminosity distances, and a Sachs equation modified by vortex/torque stress-energy. This layer supplies the observable consequences that interface with photometric catalogues, the distance ladder, and lensing statistics.  

The two layers are topologically closed: every quantity computed in the numerical layer is inserted without additional free parameters into the photometric integrals. The rapidly decaying torque belongs simultaneously to both layers—it is evaluated with numerical precision on the lattice and then propagated into the distance and magnification integrals. The continuous realisation of this photometric and distance layer, including \(\mu_{\rm lens}\), \(W_{\rm iso}\), \(\chi(z)\) and \(d_L(z)\), resides in the companion repository  
https://github.com/TheAstrographer/JCR-Kerr-Gordon_Lensing-Metric_Starter-Pack  
(with supporting discrete structures and acknowledgments also linked to https://github.com/Atrographer/Tau-the-AI).  

The framework therefore realises a single trajectory: high-redshift expansion remains essentially standard (protecting CMB/BAO consistency), while the local expansion rate is elevated by the controlled unmasking of a lattice-resident anisotropic stress. Discrepancies with high-\(z\) data would challenge the numerical core; discrepancies confined to the local \(H_0\) would challenge only the phenomenological release interpretation. The entire construction is bounded by the language of the defining documents (DISCLAIMER.tex, torque_recognition.tex, grok_acknowledgments.tex).

2. Mathematical-Numerical Layer: Exact Discrete Skeleton

The kinematic backbone is a microscopic tick lattice defined by
\[
\varepsilon=10^{-9},\qquad N_{\rm today}=10^9.
\]
The micro-comoving coordinate, scale factor and redshift are exact discrete maps:
\[
\chi_{\rm micro}=\varepsilon N,\qquad
a(N)=\exp(\chi_{\rm micro}),\qquad
z(N)=\frac{1}{a(N)}-1.
\]
Every integer tick \(N\) therefore maps onto a unique model redshift \(z(N)\) with machine precision.

Two angular quantities close an exact trigonometric identity:
\[
\psi=0.1503378808,\qquad
\Delta\phi_{\rm torque}=1.72113420759,
\]
\[
\sin(\Delta\phi_{\rm torque})=\cos\psi
\]
(to machine accuracy). This identity locks the phase relationship that governs the torque amplitude. The torque itself is fixed at the calibrated value
\[
\Delta H_{\rm torque}\big|{z=0,\,N\to N{\rm today}}=3.170\,\mathrm{km\,s^{-1}\,Mpc^{-1}},
\]
so that
\[
H_{\rm eff}(0)=73.170000
\]
exactly when the isotropic base sector sits near \(70\,\mathrm{km\,s^{-1}\,Mpc^{-1}}\).  

A rapidly decaying suppression factor (controlled by the e-folding scale \(\tau=5.8\)) leaves the high-redshift expansion essentially untouched. The same lattice that is negligible at early times reaches maximum coherence only at full maturity (\(N\to10^9\)), producing the local elevation of the expansion rate as a continuous, exponentially controlled release rather than an ad-hoc late-time correction.

3. Physical Motivation of the Torque and Screening Scales

The additive torque that appears in the Einstein-Boltzmann core,
\[
\tau(a,\chi)=\frac{3.170}{c}\cdot\frac{f_{\rm mod}(z)}{6}\cdot e^{-z/\tau_{\rm gate}}\cdot e^{-\chi/\chi_{\rm scale}},
\]
with \(\tau_{\rm gate}=5.8\), \(\chi_{\rm scale}=4500\,\mathrm{Mpc}\) and \(f_{\rm mod}(z)=1+5e^{-z/2}\), is the continuous realisation of the lattice-resident anisotropic stress.  

The damping scale \(\tau=5.8\) is anchored to cosmological transition epochs: it lies after recombination, after the main epoch of reionization, and near the matter–dark-energy equality crossover. In a vacuum-condensate / ratchet-mode picture it corresponds to the characteristic redshift at which cumulative decoherence transitions from rapid early suppression to late-time coherence buildup. The frequency-modulation scale \(\sim2\) and the amplitude-damping scale \(\sim5.8\) are deliberately separated, producing clean, non-overlapping activation regimes. Observationally the combination yields strong suppression at high \(z\) (CMB/BAO protection), moderate effects in the supernova regime, and near-full torque delivery locally.

The \(\chi\)-screening exponential further ties the torque to the growth of the comoving light-cone coordinate, guaranteeing that the anisotropic stress remains sub-dominant until the lattice reaches maturity.

4. Lensing-Photometric Layer: Observable Interface

The photometric layer converts the geometric torque into quantities that can be compared with data:

Isotropic conformal weight**  
  \[
  W_{\rm iso}(\beta,z)=1+\beta\frac{z}{1+z}\qquad(\beta=0.05).
  \]

Explicit magnification bias**  
  \[
  \mu_{\rm lens}(z)=1+0.05\frac{z}{1+z}.
  \]

Comoving and luminosity distances** obtained by direct integration of the effective expansion history  
  \[
  \chi(z)=\int_0^z\frac{c\,dz'}{H_{\rm eff}(z')},\qquad
  d_L(z)=(1+z)\chi(z).
  \]

Sachs equation** modified by the vortex / torque stress-energy contribution \(\pi_{\mu\nu}\).  

These quantities affect measured magnitudes and therefore enter the distance ladder and the supernova Hubble diagram. The layer is not an abstract numerical identity; it is the physical interface between the geometric torque and actual astronomical measurements—magnification, photometric shifts, and the recording of flux.

5. Closed Einstein-Boltzmann Core and Its Coupling to the Dual Architecture

Inside the production-grade core the state vector
\[
\mathbf{y}=(\delta_{\rm cdm},\,\theta_{\rm cdm},\,\Pi,\,\dot{\Pi},\,\chi)
\]
is evolved with respect to \(\ln a\). The effective Hubble rate that appears in every derivative is precisely the \(H_{\rm eff}\) supplied by the numerical layer (including the lattice-calibrated torque). The geometric light-cone coordinate \(\chi\) that screens the torque is the same \(\chi\) that is integrated in the photometric layer. The vortex anisotropic stress \(\Pi\) sources the gravitational slip and simultaneously contributes to the modified Sachs equation of the photometric layer.  

Thus the five-dimensional dynamical system, the discrete tick lattice, and the continuous photometric integrals form a single closed trajectory. High-redshift observables remain standard while the local expansion rate is elevated by the unmasking of the lattice-resident mechanism.

6. Interpretive Status and Bounding Language

According to the defining documents the term “Kerr–Gordon vortex” is a creative, analogical label inspired by the behaviour of a kaleidoscope magnifying glass (multiple reflections, symmetry, magnification, complex light paths). It is not claimed to be a rigorous derivation from the exact Kerr metric plus a Gordon scalar field. The construction is a phenomenological toy model / exploratory toolkit designed to explore possible resolutions of the Hubble tension through lensing-inspired corrections. All development credit belongs to Joshua Christopher Ryan.  

The late-time elevation of the expansion rate is the continuous, exponentially controlled release of anisotropic stress that has been present (but suppressed) on the lattice from early times. It is not an ad-hoc or freely adjustable late-time correction; it is the same discrete pathway reaching maximum coherence at full lattice maturity.

7. Summary Statement of the Thesis

The Production-Grade Closed Einstein-Boltzmann Core, when embedded in the dual architecture, realises a topologically closed system in which:

every discrete tick maps exactly onto a model redshift,
every phase identity closes to machine precision,
the torque amplitude is locked at \(3.170\,\mathrm{km\,s^{-1}\,Mpc^{-1}}\) and screened by both redshift and comoving distance,
the same \(H_{\rm eff}\) and \(\chi\) that govern the dynamical equations also govern the photometric distances and magnification bias,
high-\(z\) expansion is protected while the local \(H_0\) is elevated by the controlled release of lattice-resident anisotropic stress.

The mathematical-numerical layer supplies the exact clockwork; the lensing-photometric layer supplies the observable face. Together they constitute a single, gated, and bounded trajectory from the microscopic lattice to the measured distance ladder.
